\section{Conclusion and future work}
\label{sec:conclusion}

This work tackles EEG denoising under the cross-subject generalization challenge that limits multi-subject deep-learning pipelines. We designed and implemented \saddpm{}, a framework that couples a denoising diffusion probabilistic model with a subject-aware mechanism: trainable low-dimensional subject embeddings are injected as conditioning signals at multiple network levels, so the model adapts its reverse denoising trajectory to each individual while retaining the training stability of diffusion. For target subjects without a pre-trained embedding, a transductive, label-free step estimates only the embedding from unlabeled target-domain EEG, keeping the shared backbone fixed.

Across two complementary protocols on BCI Competition~IV-2a, the approach proves viable. Under downstream motor-imagery classification, \saddpm{} is more stable than the classical ICA baseline, with a mean-accuracy gain concentrated in cross-subject transfer and generated signals that preserve subject-specific structure. Under the EEGdenoiseNet paired ground-truth protocol, the conditional variant \saddpmcond{} is competitive with supervised CNN denoisers on ocular artifacts, attains high-fidelity multi-channel denoising (CC~$\approx0.99$), and exhibits subject conditioning that is functionally relevant precisely when artifacts are subject-specific. We read these as evidence of feasibility rather than superiority over the strongest learned denoisers.

Current subject adaptation still relies on unlabeled target data; reducing this dependence through transfer, few-shot, or zero-shot strategies---together with stronger baselines, more datasets and artifact types, end-to-end joint denoising--decoding, and faster sampling for low-latency BCI---are natural next steps. More broadly, \saddpm{} offers a novel and practical route to subject-adaptive generative denoising for multi-subject EEG, extends diffusion models to individual-aware neural time-series processing, and frames EEG
denoising as an autonomous and adaptive systems problem: a learning-based BCI component that must stay reliable under changing users and cross-user
distribution shift~\cite{casimiro2024selfadapting}.